\documentclass[11pt,oneside]{article}

\usepackage[T1]{fontenc}
\usepackage{lmodern}
\usepackage{microtype}
\usepackage[margin=1in,headheight=14pt]{geometry}
\usepackage{amsmath,amssymb,amsthm,mathtools}
\usepackage{booktabs,tabularx,array,longtable}
\usepackage{enumitem}
\usepackage{xcolor}
\usepackage{xurl}
\usepackage{hyperref}
\usepackage{fancyhdr}
\usepackage{titlesec}

\definecolor{linkblue}{RGB}{25,82,140}
\definecolor{proofgreen}{RGB}{17,105,72}
\definecolor{warningred}{RGB}{145,45,45}
\definecolor{softgray}{RGB}{245,246,248}
\definecolor{rulegray}{RGB}{105,110,118}

\hypersetup{
  colorlinks=true,
  linkcolor=linkblue,
  citecolor=proofgreen,
  urlcolor=linkblue,
  pdftitle={Additive Moment Rigidity and the Alaoglu--Erdos Problem},
  pdfauthor={Independent research continuation},
  pdfsubject={Hausdorff moments, exact denominator defects, rational quadrature, and Muntz incompleteness}
}
\urlstyle{same}

\pagestyle{fancy}
\fancyhf{}
\lhead{Alaoglu--Erd\H{o}s: additive moment continuation}
\rhead{August 21, 2026}
\cfoot{\thepage}
\renewcommand{\headrulewidth}{0.4pt}

\titleformat{\section}{\Large\bfseries}{\thesection}{0.75em}{}
\titleformat{\subsection}{\large\bfseries}{\thesubsection}{0.75em}{}
\titleformat{\subsubsection}{\normalsize\bfseries}{\thesubsubsection}{0.75em}{}

\setlength{\parindent}{0pt}
\setlength{\parskip}{0.55em}
\setlength{\emergencystretch}{3em}
\setlist[itemize]{leftmargin=2em,itemsep=0.25em,topsep=0.35em}
\setlist[enumerate]{leftmargin=2.2em,itemsep=0.25em,topsep=0.35em}

\newtheorem{theorem}{Theorem}[section]
\newtheorem{proposition}[theorem]{Proposition}
\newtheorem{lemma}[theorem]{Lemma}
\newtheorem{corollary}[theorem]{Corollary}
\newtheorem{conjecture}[theorem]{Conjecture}
\newtheorem{target}[theorem]{Research target}
\theoremstyle{definition}
\newtheorem{definition}[theorem]{Definition}
\newtheorem{observation}[theorem]{Observation}
\newtheorem{warning}[theorem]{Warning}
\theoremstyle{remark}
\newtheorem{remark}[theorem]{Remark}

\newcommand{\N}{\mathbb N}
\newcommand{\Nzero}{\mathbb N_0}
\newcommand{\Z}{\mathbb Z}
\newcommand{\Q}{\mathbb Q}
\newcommand{\R}{\mathbb R}
\newcommand{\C}{\mathbb C}
\newcommand{\Ssmooth}{\mathcal S_{2,3}}
\newcommand{\dd}{\,\mathrm d}
\newcommand{\supp}{\operatorname{supp}}
\newcommand{\Span}{\operatorname{span}}
\newcommand{\cl}{\operatorname{cl}}
\newcommand{\AEname}{Alaoglu--Erd\H{o}s}
\newcommand{\FE}{Four Exponentials Conjecture}

\newenvironment{statusbox}
{\begin{center}\begin{minipage}{0.92\textwidth}\small\hrule\vspace{0.45em}}
{\vspace{0.35em}\hrule\end{minipage}\end{center}}

\begin{document}
\hypersetup{pageanchor=false}
\pagenumbering{gobble}

\begin{titlepage}
\centering
\vspace*{1.0cm}
{\Huge\bfseries Additive Moment Rigidity\par}
\vspace{0.18cm}
{\Huge\bfseries and the Alaoglu--Erd\H{o}s Problem\par}
\vspace{0.65cm}
{\Large A Hausdorff--M\"untz continuation beyond the supplied unified report\par}
\vspace{1.0cm}
{\large Independent research report\par}
\vspace{0.15cm}
{\large August 21, 2026\par}
\vfill
\begin{minipage}{0.91\textwidth}
\small
\textbf{Status statement.}
The implication
\[
  2^x,3^x\in\Z \quad\Longrightarrow\quad x\in\Z
\]
is not proved in this report.  As of the date above it remains open, and a positive answer
follows from the Four Exponentials Conjecture.  This document nevertheless gives a detailed,
nonduplicative investigation of a new route based on the additive relation $3=2+1$.
All theorems labeled as such below are proved.  The main outcome is both constructive and
negative: a hypothetical counterexample produces an explicit Hausdorff moment sequence with
several rational moments and exact prime-power denominators, but finite rational quadrature and
a M\"untz incompleteness theorem show rigorously why linear moment positivity cannot finish the
problem.  The remaining step is isolated as a nonlinear arithmetic-rigidity target.
\end{minipage}
\vfill
\end{titlepage}
\hypersetup{pageanchor=true}
\pagenumbering{roman}

\begin{abstract}
Let $x\in\R$ satisfy $2^x,3^x\in\Z$, and suppose $x$ is not an integer.  Writing
$x=n+\alpha$ with $0<\alpha<1$ gives rational numbers
$u=2^\alpha\in\Z[1/2]$ and $v=3^\alpha\in\Z[1/3]$.  The supplied unified report already
pushes transcendence, determinant, $p$-adic, matrix, probability, symbolic-dynamical, and
polynomial-approximation methods to their current barriers.  The present report instead
exploits $3=2+1$ by studying the increments
\[
 d_j=(j+2)^\alpha-(j+1)^\alpha.
\]
They are the moments of the explicit positive measure
\[
 \dd\mu_\alpha(q)=\frac{\alpha}{\Gamma(1-\alpha)}
 (1-q)(-\log q)^{-\alpha-1}\dd q,
 \qquad 0<q<1.
\]
Consequently $(d_j)$ is strictly completely monotone, all its Hankel matrices are positive
definite, and it is strictly log-convex.  The consecutive $3$-smooth pairs are exactly
$(1,2),(2,3),(3,4),(8,9)$, so a counterexample makes precisely
$d_0,d_1,d_2,d_7$ immediately rational.  We compute exact reduced denominators for the first
Hankel defect and three higher log-convexity defects; the numerators are positive odd integers and are also prime to $3$ whenever a triadic denominator is present.

The decisive audit is then twofold.  First, every finite set of these rational moment
constraints has a positive atomic realization with rational nodes and rational weights; an
explicit two-atom formula already reproduces $d_0,d_1,d_2$.  Second, the exponents
$2^a3^b-1$ have a convergent reciprocal sum.  The full M\"untz theorem therefore supplies
bounded, compactly supported perturbations of $\mu_\alpha$ that preserve \emph{all} linear
$3$-smooth observations
\[
 \int_0^1\frac{1-q^{N-1}}{1-q}\dd\mu(q),\qquad N=2^a3^b,
\]
while leaving the endpoint singularities and moment asymptotics unchanged.  Thus no proof can
come from linear positivity, finite Hankel certificates, measure reconstruction, or endpoint
tail data alone.  A successful argument must use nonlinear coupling between the explicit
one-parameter density, multiplicativity, and the separated dyadic/triadic denominator
arithmetic.  That exact target is formulated at the end.
\end{abstract}

\tableofcontents
\clearpage
\pagenumbering{arabic}

\section{Scope, current status, and novelty}

\subsection{The problem}

The \AEname{} integer-exponent conjecture is the following elementary-looking statement.

\begin{conjecture}[\AEname{}]
If $x\in\R$ and $2^x,3^x\in\Z$, then $x\in\Z$.
\end{conjecture}

The question appears in the circle of problems originating in Alaoglu and Erd\H{o}s's work on
highly composite and colossally abundant numbers \cite{AlaogluErdos1944}.  Waldschmidt's
modern survey states the question explicitly and records that a positive answer follows from
the four exponentials problem; the latter remains open \cite{Waldschmidt2023}.  Recent work on
matrices of logarithms likewise formulates, rather than resolves, a matrix-coefficient
conjecture designed to control vanishing logarithmic determinants \cite{DasguptaKakde2024}.

The statement is easy for rational or algebraic exponents.  Its difficulty is that any
counterexample must be transcendental and creates a singular $2\times2$ matrix of logarithms
at exactly the size not reached by the Six Exponentials Theorem.  This is not a reason to stop;
it is a constraint on what a valid proof must add.

\subsection{Relation to the supplied unified report}

The supplied source
\path{Article/alaoglu-erdos-unified-report/alaoglu-erdos-unified-report.tex}
is a very large audit of the existing project.  It already develops:

\begin{itemize}
\item the reduction to Four Exponentials and the Gelfond--Schneider consequences;
\item exact conditional free-monoid structure if one counterexample exists;
\item ordinary and generalized Vandermonde determinants, interpolation, Loewner matrices,
      HCIZ asymptotics, and finite differences;
\item $p$-adic valuation cascades, tropical assignment layers, and formal Lean modules;
\item rational-point counting, Sturmian and automatic-sequence models, probability,
      information-theoretic, Mahler, spectral, and polynomial-approximation diagnostics.
\end{itemize}

This report does not repeat those developments.  It begins only with the minimal reductions
needed for a standalone argument, and then studies a direction absent from the supplied text:
\emph{the Hausdorff moment structure forced by the additive adjacency of the $3$-smooth
numbers $1,2,3,4,8,9$}.  The strongest conclusions are new exact denominator theorems and two
no-go theorems---finite rational quadrature and infinite M\"untz invisibility.

\begin{statusbox}
\textbf{Trust boundary.}  The report proves every displayed theorem by ordinary mathematical
arguments.  It does not claim a Lean formalization.  Classical inputs used as black boxes are
Gelfond--Schneider, Catalan--Mih\u{a}ilescu, and the full M\"untz theorem in $L^1$ on an interval
away from the origin.
\end{statusbox}

\section{Minimal reductions and the four-exponentials wall}

\begin{proposition}[Elementary and transcendence reductions]\label{prop:minimal-reductions}
Suppose $2^x,3^x\in\Z$.
\begin{enumerate}[label=\textup{(\roman*)}]
\item $x\ge0$.
\item If $x\in\Q$, then $x\in\Nzero$.
\item If $x$ is algebraic, then $x\in\Nzero$.
\item Assuming the Four Exponentials Conjecture, $x\in\Nzero$.
\end{enumerate}
Consequently every hypothetical counterexample is a positive transcendental irrational.
\end{proposition}

\begin{proof}
If $x<0$, then $0<2^x<1$, so $2^x$ cannot be an integer.  This proves (i).

For (ii), write $x=a/b$ in lowest terms with $a\ge0$ and $b\ge1$.  If $M=2^x\in\N$, then
$M^b=2^a$.  Taking the $2$-adic valuation gives $b\,v_2(M)=a$.  Hence $b\mid a$; coprimality
forces $b=1$.

For (iii), if $x$ were algebraic and irrational, the Gelfond--Schneider theorem would make
$2^x$ transcendental, contradicting $2^x\in\Z$.  Part (ii) handles the remaining algebraic
case.

For (iv), assume $x$ is not rational.  The pair $(\log2,\log3)$ is linearly independent over
$\Q$, because $2^a3^b=1$ with $a,b\in\Z$ forces $a=b=0$.  The pair $(1,x)$ is also
$\Q$-linearly independent.  The four numbers
\[
 e^{(\log2)\cdot1}=2,\quad e^{(\log2)x}=2^x,\quad
 e^{(\log3)\cdot1}=3,\quad e^{(\log3)x}=3^x
\]
are algebraic.  Four Exponentials would force one to be transcendental, a contradiction.
Thus $x\in\Q$, and (ii) applies.
\end{proof}

Equivalently, a counterexample gives positive integers $M,A$ with
\begin{equation}\label{eq:log-det}
 \frac{\log M}{\log2}=\frac{\log A}{\log3}=x,
 \qquad
 \det\begin{pmatrix}\log2&\log3\\ \log M&\log A\end{pmatrix}=0.
\end{equation}
The row and column pairs are $\Q$-independent when $x$ is irrational.  Equation
\eqref{eq:log-det} is therefore a concrete $2\times2$ instance of the four-exponentials
barrier.  The rest of this report asks whether the special additive relation $3=2+1$ adds
something that the logarithmic determinant does not see.

\section{Fractional normalization and prime-separated denominators}

Let $x$ be a hypothetical counterexample.  Put
\begin{equation}\label{eq:frac-normalization}
 n=\lfloor x\rfloor,\qquad \alpha=x-n\in(0,1),\qquad
 u=2^\alpha,\qquad v=3^\alpha.
\end{equation}
Since $2^x$ and $3^x$ are integers,
\[
 u=\frac{2^x}{2^n}\in\Q,
 \qquad
 v=\frac{3^x}{3^n}\in\Q.
\]
Write these fractions in lowest terms as
\begin{equation}\label{eq:uv-reduced}
 u=\frac{P}{2^r},\qquad v=\frac{Q}{3^s},
\end{equation}
where
\begin{equation}\label{eq:rs-conditions}
 1\le r\le n,
 \qquad 0\le s\le n,
 \qquad P\ \text{is odd},
 \qquad 3\nmid Q.
\end{equation}
The inequality $r\ge1$ follows from $1<u<2$: a rational integer strictly between $1$ and $2$
does not exist.  The case $s=0$ is not excluded a priori, because the integer $v=2$ lies in
$(1,3)$.

Define the $3$-smooth semigroup
\[
 \Ssmooth=\{2^a3^b:a,b\in\Nzero\}.
\]
For every $N=2^a3^b\in\Ssmooth$,
\begin{equation}\label{eq:s-smooth-rationality}
 N^\alpha=u^av^b\in\Q.
\end{equation}
Thus one hypothetical solution gives infinitely many rational values of the concave function
$t\mapsto t^\alpha$---but only on a logarithmically sparse multiplicative semigroup.

Two elementary endpoint gaps will be useful for calibration.

\begin{lemma}[The one-coordinate endpoint gaps]\label{lem:trivial-gaps}
With the notation above,
\begin{equation}\label{eq:trivial-gaps}
 \alpha\ge2^{-r},
 \qquad
 1-\alpha\ge\frac{2^{-r}}{2\log2}.
\end{equation}
In particular both distances are bounded below exponentially in $n$.
\end{lemma}

\begin{proof}
The reduced numerator of $u-1=(P-2^r)/2^r$ is a positive odd integer, so
$u-1\ge2^{-r}$.  Convexity of $2^t$ on $[0,1]$ places its graph below the chord joining
$(0,1)$ and $(1,2)$; hence $2^\alpha-1\le\alpha$.  This proves the first inequality.

Similarly, $2-u=(2^{r+1}-P)/2^r\ge2^{-r}$.  Also
\[
 2-2^\alpha=\int_\alpha^1 2^t\log2\dd t
 \le2\log2\,(1-\alpha).
\]
Combining the two bounds proves the second inequality.
\end{proof}

These simple gaps are important when evaluating later determinant bounds: a new positivity
defect is only useful if it beats information already visible in the single dyadic coordinate.

\section{The Hausdorff moment representation}

\subsection{The increment sequence}

For $j\in\Nzero$ define
\begin{equation}\label{eq:dj-def}
 d_j=(j+2)^\alpha-(j+1)^\alpha.
\end{equation}
Because $0<\alpha<1$, the sequence is positive and strictly decreasing.  More is true: it is a
strict Hausdorff moment sequence with an explicit density.

\begin{theorem}[Explicit Hausdorff measure]\label{thm:hausdorff}
For $0<\alpha<1$, define a measure on $(0,1)$ by
\begin{equation}\label{eq:mu-density}
 \dd\mu_\alpha(q)=w_\alpha(q)\dd q,
 \qquad
 w_\alpha(q)=\frac{\alpha}{\Gamma(1-\alpha)}
 (1-q)(-\log q)^{-\alpha-1}.
\end{equation}
Then $\mu_\alpha$ is finite, positive, and has full support in $[0,1]$.  Moreover
\begin{equation}\label{eq:moment-representation}
 d_j=\int_0^1q^j\dd\mu_\alpha(q)
 \qquad(j\in\Nzero).
\end{equation}
More generally, for every real $z>-1$,
\begin{equation}\label{eq:mellin-extension}
 \int_0^1q^z\dd\mu_\alpha(q)
 =(z+2)^\alpha-(z+1)^\alpha.
\end{equation}
\end{theorem}

\begin{proof}
For $t>0$, the gamma integral gives
\[
 t^{\alpha-1}
 =\frac1{\Gamma(1-\alpha)}
 \int_0^\infty e^{-ts}s^{-\alpha}\dd s.
\]
Since
\[
 d_j=\int_{j+1}^{j+2}\alpha t^{\alpha-1}\dd t,
\]
Tonelli's theorem yields
\begin{align*}
 d_j
 &=\frac{\alpha}{\Gamma(1-\alpha)}
   \int_0^\infty s^{-\alpha}
   \left(\int_{j+1}^{j+2}e^{-ts}\dd t\right)\dd s\\
 &=\frac{\alpha}{\Gamma(1-\alpha)}
   \int_0^\infty e^{-js}e^{-s}(1-e^{-s})s^{-\alpha-1}\dd s.
\end{align*}
Put $q=e^{-s}$.  Then $\dd s=-\dd q/q$, and the last expression becomes
\[
 \frac{\alpha}{\Gamma(1-\alpha)}
 \int_0^1q^j(1-q)(-\log q)^{-\alpha-1}\dd q,
\]
which is \eqref{eq:moment-representation}.

Near $q=1$, $-\log q\sim1-q$, so the density is asymptotic to a positive constant times
$(1-q)^{-\alpha}$ and is integrable because $\alpha<1$.  Near $q=0$ the change of variables
$q=e^{-s}$ shows integrability.  The density is strictly positive on $(0,1)$, hence the support
is all of $[0,1]$.  Replacing the integer $j$ by any real $z>-1$ in the same computation gives
\eqref{eq:mellin-extension}.
\end{proof}

\subsection{Complete monotonicity, Hankel positivity, and log-convexity}

\begin{corollary}[Strict complete monotonicity]\label{cor:complete-monotone}
For all $j,k\in\Nzero$,
\begin{equation}\label{eq:complete-monotone}
 (-1)^k\Delta^kd_j
 =\int_0^1q^j(1-q)^k\dd\mu_\alpha(q)>0,
\end{equation}
where $\Delta d_j=d_{j+1}-d_j$.
\end{corollary}

\begin{proof}
Apply $\Delta^k$ to $q^j$ inside \eqref{eq:moment-representation}.  Since
$\Delta^kq^j=q^j(q-1)^k$, the formula follows.  Strict positivity uses the positive density on
$(0,1)$.
\end{proof}

\begin{corollary}[Strict Hankel positivity]\label{cor:hankel}
For every $m\ge0$, the Hankel matrix
\[
 H_m=(d_{i+j})_{0\le i,j\le m}
\]
is positive definite.  In particular,
\begin{equation}\label{eq:first-hankel}
 d_0d_2-d_1^2>0.
\end{equation}
\end{corollary}

\begin{proof}
For a nonzero real vector $(c_0,\ldots,c_m)$,
\[
 \sum_{i,j=0}^m c_ic_jd_{i+j}
 =\int_0^1\left(\sum_{i=0}^mc_iq^i\right)^2\dd\mu_\alpha(q)>0.
\]
A nonzero polynomial cannot vanish on the full support of $\mu_\alpha$.
\end{proof}

\begin{corollary}[Strict log-convexity]\label{cor:logconvex}
If $0\le i<j<k$, then
\begin{equation}\label{eq:logconvex-general}
 d_j^{\,k-i}<d_i^{\,k-j}d_k^{\,j-i}.
\end{equation}
\end{corollary}

\begin{proof}
Write
\[
 j=\frac{k-j}{k-i}i+\frac{j-i}{k-i}k
\]
and apply H\"older's inequality to
$q^j=(q^i)^{(k-j)/(k-i)}(q^k)^{(j-i)/(k-i)}$.  Equality would force $q^i$ and $q^k$ to be
proportional almost everywhere, impossible on a measure with full support.
\end{proof}

\subsection{Tail behavior}

The mean-value theorem gives
\begin{equation}\label{eq:tail}
 d_j=\alpha\xi_j^{\alpha-1}
 \quad\text{for some }\xi_j\in(j+1,j+2),
 \qquad
 d_j\sim\alpha j^{\alpha-1}.
\end{equation}
Thus the moment decay is polynomial, not exponential.  In measure language, this is caused by
the integrable singularity of $w_\alpha$ at $q=1$.  This observation will separate the genuine
infinite moment sequence from the finite atomic models constructed later.

\section{Exactly which adjacent \texorpdfstring{$3$}{3}-smooth values are rational?}

The formula $d_j=(j+2)^\alpha-(j+1)^\alpha$ is rational whenever both consecutive integers
$j+1$ and $j+2$ are $3$-smooth.  Those pairs can be classified exactly.

\begin{theorem}[Consecutive $3$-smooth pairs]\label{thm:consecutive-smooth}
The positive consecutive pairs $(N,N+1)$ for which both entries belong to $\Ssmooth$ are
\begin{equation}\label{eq:consecutive-list}
 (1,2),\qquad(2,3),\qquad(3,4),\qquad(8,9).
\end{equation}
\end{theorem}

\begin{proof}
The pair $(1,2)$ is immediate.  Suppose $N>1$.  Since $\gcd(N,N+1)=1$ and both numbers have
all prime factors in $\{2,3\}$, their prime supports are disjoint.  Hence one is a pure power
of $2$ and the other a pure power of $3$:
\[
 \{N,N+1\}=\{2^a,3^b\}
 \qquad(a,b\ge1).
\]
If $a,b>1$, Catalan--Mih\u{a}ilescu says that the unique pair of nontrivial perfect powers that
differ by $1$ is $8$ and $9$ \cite{Mihailescu2004}.  If one exponent equals $1$, direct
inspection gives $2,3$ and $3,4$.  No other cases occur.
\end{proof}

\begin{corollary}[The four immediately rational moments]\label{cor:four-moments}
Under the hypothetical counterexample normalization \eqref{eq:frac-normalization},
\begin{equation}\label{eq:four-moments}
 \begin{aligned}
 d_0&=u-1,\\
 d_1&=v-u,\\
 d_2&=u^2-v,\\
 d_7&=v^2-u^3,
 \end{aligned}
\end{equation}
so $d_0,d_1,d_2,d_7\in\Q_{>0}$.
\end{corollary}

The list is exact for unit increments.  Of course every $3$-smooth value gives a rational
\emph{partial sum} of the moments:
\begin{equation}\label{eq:partial-sums}
 N^\alpha-1
 =\sum_{j=0}^{N-2}d_j
 =\int_0^1\frac{1-q^{N-1}}{1-q}\dd\mu_\alpha(q),
 \qquad N\in\Ssmooth,
\end{equation}
but only the four pairs in \eqref{eq:consecutive-list} isolate single moments directly.

\section{Exact prime-power denominators of moment defects}

\subsection{Arithmetic normalization}

Retain \eqref{eq:uv-reduced}.  Define the positive integers
\begin{equation}\label{eq:A-defs}
 \begin{aligned}
 A_0&=P-2^r,\\
 A_1&=2^rQ-3^sP,\\
 A_2&=3^sP^2-2^{2r}Q,\\
 A_7&=2^{3r}Q^2-3^{2s}P^3.
 \end{aligned}
\end{equation}
Then
\begin{equation}\label{eq:dA}
 d_0=\frac{A_0}{2^r},\qquad
 d_1=\frac{A_1}{2^r3^s},\qquad
 d_2=\frac{A_2}{2^{2r}3^s},\qquad
 d_7=\frac{A_7}{2^{3r}3^{2s}}.
\end{equation}
Every $A_i$ is odd.  If $s\ge1$, then
\begin{equation}\label{eq:A-mod3}
 3\nmid A_1A_2A_7.
\end{equation}
Indeed, the four numerators in \eqref{eq:A-defs} are respectively odd-minus-even,
even-minus-odd, odd-minus-even, and even-minus-odd.  Modulo $3$, when $s\ge1$,
\[
 A_1\equiv2^rQ,
 \quad A_2\equiv-2^{2r}Q,
 \quad A_7\equiv2^{3r}Q^2,
\]
all nonzero.

\subsection{The first Hankel defect}

\begin{theorem}[Exact denominator of the first Hankel defect]\label{thm:H-denom}
Let
\begin{equation}\label{eq:H-def}
 \mathcal H=d_0d_2-d_1^2.
\end{equation}
Then
\begin{equation}\label{eq:H-uv}
 \mathcal H
 =(u-1)(u^2-v)-(v-u)^2
 =u^3-2u^2+uv+v-v^2>0,
\end{equation}
and in lowest terms
\begin{equation}\label{eq:H-exact}
 \mathcal H
 =\frac{3^sA_0A_2-2^rA_1^2}{2^{3r}3^{2s}}.
\end{equation}
In particular, the reduced denominator is exactly $2^{3r}3^{2s}$ and
\begin{equation}\label{eq:H-lower}
 \mathcal H\ge2^{-3r}3^{-2s}.
\end{equation}
\end{theorem}

\begin{proof}
Positivity is Corollary~\ref{cor:hankel}; the expression in $u,v$ is a direct expansion.
Substitution of \eqref{eq:dA} gives \eqref{eq:H-exact}.  Its numerator is odd because the
first term is odd and the second is even.  If $s\ge1$, the numerator is congruent modulo $3$
to $-2^rA_1^2\ne0$.  If $s=0$, no factor of $3$ occurs in the denominator.  Hence no factor
$2$ or $3$ cancels, proving exactness.  The positive numerator is at least $1$.
\end{proof}

The double-integral identity
\begin{equation}\label{eq:H-variance}
 \mathcal H
 =\frac12\int_0^1\int_0^1(q-t)^2\dd\mu_\alpha(q)\dd\mu_\alpha(t)
\end{equation}
shows that $\mathcal H$ is the unnormalized variance of the Hausdorff measure.  Theorem
\ref{thm:H-denom} therefore turns a continuous variance into a positive integer divided by a
fully determined prime-power denominator.

\subsection{Three higher log-convexity defects}

Apply Corollary~\ref{cor:logconvex} to the triples $(0,1,7)$, $(1,2,7)$, and $(0,2,7)$.

\begin{theorem}[Exact higher defects]\label{thm:higher-defects}
The rational numbers
\begin{align}
 \mathcal F_{017}&=d_0^6d_7-d_1^7,\label{eq:F017}\\
 \mathcal F_{127}&=d_1^5d_7-d_2^6,\label{eq:F127}\\
 \mathcal F_{027}&=d_0^5d_7^2-d_2^7\label{eq:F027}
\end{align}
are strictly positive.  Their exact reduced forms are
\begin{align}
 \mathcal F_{017}
 &=\frac{3^{5s}A_0^6A_7-2^{2r}A_1^7}
 {2^{9r}3^{7s}},\label{eq:F017-exact}\\
 \mathcal F_{127}
 &=\frac{2^{4r}A_1^5A_7-3^sA_2^6}
 {2^{12r}3^{7s}},\label{eq:F127-exact}\\
 \mathcal F_{027}
 &=\frac{2^{3r}3^{3s}A_0^5A_7^2-A_2^7}
 {2^{14r}3^{7s}}.\label{eq:F027-exact}
\end{align}
Each displayed numerator is a positive odd integer; if $s\ge1$, it is coprime to $6$.  In every case it is coprime to the displayed denominator.  Consequently
\begin{equation}\label{eq:F-lowers}
 \mathcal F_{017}\ge2^{-9r}3^{-7s},\quad
 \mathcal F_{127}\ge2^{-12r}3^{-7s},\quad
 \mathcal F_{027}\ge2^{-14r}3^{-7s}.
\end{equation}
\end{theorem}

\begin{proof}
Strict positivity is exactly \eqref{eq:logconvex-general} for the three indicated index
triples.  Equations \eqref{eq:F017-exact}--\eqref{eq:F027-exact} follow by substituting
\eqref{eq:dA} and taking a common denominator.

For \eqref{eq:F017-exact}, the first numerator term is odd and the second even, so the
numerator is odd.  If $s\ge1$, it is congruent to $-2^{2r}A_1^7\ne0\pmod3$.
For \eqref{eq:F127-exact}, the first term is even and the second odd, and modulo $3$ the
numerator is $2^{4r}A_1^5A_7\ne0$.  For \eqref{eq:F027-exact}, the first term is even and the
second odd, and modulo $3$ the numerator is $-A_2^7\ne0$.  The case $s=0$ again has no
triadic denominator.  Exactness and the lower bounds follow.
\end{proof}

\begin{table}[htbp]
\centering
\caption{Exact arithmetic defects supplied by the four rational moments.}
\label{tab:defects}
\small
\begin{tabularx}{\textwidth}{@{}l l l X@{}}
\toprule
Defect & Positivity source & Exact denominator & Numerator information\\
\midrule
$d_0d_2-d_1^2$
& $2\times2$ Hankel
& $2^{3r}3^{2s}$
& positive; coprime to the denominator\\
$d_0^6d_7-d_1^7$
& log-convex $(0,1,7)$
& $2^{9r}3^{7s}$
& positive; coprime to the denominator\\
$d_1^5d_7-d_2^6$
& log-convex $(1,2,7)$
& $2^{12r}3^{7s}$
& positive; coprime to the denominator\\
$d_0^5d_7^2-d_2^7$
& log-convex $(0,2,7)$
& $2^{14r}3^{7s}$
& positive; coprime to the denominator\\
\bottomrule
\end{tabularx}
\end{table}

\subsection{What the exact defects do and do not buy}

The exact denominator is genuine arithmetic information, but its immediate size consequence is
not enough.  For example, since $d_7<d_0$ and $d_0=2^\alpha-1\le\alpha$,
\[
 0<\mathcal F_{017}<d_0^7\le\alpha^7.
\]
Combining this with \eqref{eq:F-lowers} gives
\begin{equation}\label{eq:F-alpha-gap}
 \alpha\ge2^{-9r/7}3^{-s}.
\end{equation}
This is weaker than the elementary dyadic estimate $\alpha\ge2^{-r}$ in
Lemma~\ref{lem:trivial-gaps} for many parameter ranges, and in any case remains exponentially
small in $n$.  Similarly, an upper bound for $\mathcal H$ near $\alpha=1$ yields only another
exponential endpoint gap.  The defects therefore provide exact certificates and useful
formalization targets, but not a contradiction.

This failure motivates the next question: can one combine many moment inequalities?  The
answer is unexpectedly rigid in the opposite direction.  Every finite collection can be
simulated by a rational atomic measure.

\section{Finite rational quadrature: a no-go theorem}

\subsection{An explicit two-atom model}

\begin{theorem}[Rational two-atom reproduction]\label{thm:two-atom}
Let $m_0,m_1,m_2\in\Q_{>0}$ satisfy
\[
 m_0>m_1>m_2,
 \qquad m_0m_2>m_1^2.
\]
Then the positive atomic measure
\begin{equation}\label{eq:two-atom-measure}
 \nu=
 \left(m_0-\frac{m_1^2}{m_2}\right)\delta_0
 +\frac{m_1^2}{m_2}\,\delta_{m_2/m_1}
\end{equation}
has rational nodes and rational weights and satisfies
\begin{equation}\label{eq:two-atom-moments}
 \int q^j\dd\nu(q)=m_j
 \qquad(j=0,1,2).
\end{equation}
In particular, under a hypothetical \AEname{} counterexample, the first three moments
$d_0,d_1,d_2$ have a completely rational positive model unrelated to the transcendental
parameter $\alpha$.
\end{theorem}

\begin{proof}
The node $t=m_2/m_1$ lies in $(0,1)$.  The weight at $0$ is
$(m_0m_2-m_1^2)/m_2>0$, and the other weight is positive.  Direct calculation gives
\[
 \nu([0,1])=m_0,
 \quad
 \int q\dd\nu=\frac{m_1^2}{m_2}\frac{m_2}{m_1}=m_1,
 \quad
 \int q^2\dd\nu=\frac{m_1^2}{m_2}\frac{m_2^2}{m_1^2}=m_2.
\]
All quantities are rational.
\end{proof}

For $m_j=d_j$, the positivity of the first weight is exactly the Hankel inequality
$d_0d_2-d_1^2>0$.  Thus the first nontrivial Hankel certificate is not merely compatible with
rational data; it canonically constructs a rational quadrature formula for it.

\subsection{A general rational finite-observable theorem}

The preceding phenomenon is not confined to three moments.

\begin{lemma}[Rational interior quadrature]\label{lem:rational-quadrature}
Let $f_1,\ldots,f_m\in\Q[q]$ be linearly independent real polynomials, and let $\mu$ be a
finite positive measure with a strictly positive density on $(0,1)$.  Suppose
\[
 y_i=\int_0^1 f_i(q)\dd\mu(q)\in\Q
 \qquad(1\le i\le m).
\]
Then there are finitely many rational nodes $q_1,\ldots,q_L\in(0,1)$ and positive rational
weights $w_1,\ldots,w_L$ such that
\begin{equation}\label{eq:rational-quadrature-general}
 y_i=\sum_{\ell=1}^Lw_\ell f_i(q_\ell)
 \qquad(1\le i\le m).
\end{equation}
\end{lemma}

\begin{proof}
Set $F(q)=(f_1(q),\ldots,f_m(q))\in\R^m$, and let $C$ be the convex cone generated by
$F((0,1))$.  The vector $y=(y_1,\ldots,y_m)$ lies in the interior of $C$.  Indeed, if a
nonzero linear functional $c\in(\R^m)^*$ supported $C$ at $y$, then the polynomial
$p(q)=c\cdot F(q)$ would be nonnegative on $(0,1)$ and satisfy
$\int p\dd\mu=0$.  Positivity of the density would force $p$ to vanish identically, contrary
to linear independence of the $f_i$.

By finite-dimensional convex geometry, an interior point of a generated cone lies in the
interior of a cone generated by finitely many of its generating rays.  Choose such nodes and
then perturb them slightly to distinct rational nodes.  Interior containment is open, so it
persists under a sufficiently small perturbation.  We obtain a rational matrix
$B=(f_i(q_\ell))$ for which the system $Bw=y$ has a strictly positive real solution.  The
affine solution space is defined over $\Q$; its rational points are dense.  Hence it has a
strictly positive rational solution, which gives \eqref{eq:rational-quadrature-general}.
\end{proof}

\begin{corollary}[Finite simulation of all available additive data]\label{cor:finite-simulation}
Assume a hypothetical counterexample.  Then:
\begin{enumerate}[label=\textup{(\roman*)}]
\item the four rational moments $d_0,d_1,d_2,d_7$ are reproduced exactly by a finite positive
      measure with rational atoms and rational weights;
\item for every finite set $N_1,\ldots,N_m\in\Ssmooth$, there is such a rational atomic
      measure $\nu$ satisfying
\[
 \int_0^1\frac{1-q^{N_i-1}}{1-q}\dd\nu(q)=N_i^\alpha-1
 \qquad(1\le i\le m).
\]
\end{enumerate}
\end{corollary}

\begin{proof}
For (i), use Lemma~\ref{lem:rational-quadrature} with
$f(q)=1,q,q^2,q^7$ and \eqref{eq:four-moments}.  For (ii), use the linearly independent
polynomials
\[
 K_{N_i-1}(q)=\frac{1-q^{N_i-1}}{1-q}=1+q+\cdots+q^{N_i-2}
\]
after deleting repetitions.  Their integrals are rational by \eqref{eq:partial-sums} and
\eqref{eq:s-smooth-rationality}.
\end{proof}

\begin{warning}[Finite positivity cannot prove the conjecture]
Any argument that uses only a finite number of linear moment identities and positivity of
integrals against nonnegative polynomials can be reproduced inside a completely rational
finite atomic model.  Therefore no finite Hankel matrix, no finite list of complete-monotonicity
inequalities, and no finite positive-kernel certificate can distinguish the putative
counterexample from rational quadrature data.  A proof must be genuinely infinite or use
additional nonlinear arithmetic.
\end{warning}

\section{The complete \texorpdfstring{$3$}{3}-smooth observation family}

\subsection{The observation kernels}

For $m\ge1$, put
\begin{equation}\label{eq:Km}
 K_m(q)=\frac{1-q^m}{1-q}=1+q+\cdots+q^{m-1}.
\end{equation}
Equation \eqref{eq:partial-sums} becomes
\begin{equation}\label{eq:Km-data}
 \int_0^1K_{N-1}(q)\dd\mu_\alpha(q)=N^\alpha-1
 \qquad(N\in\Ssmooth,\ N\ge2).
\end{equation}
Under a counterexample, the right side is not merely rational; it factors coherently as
$u^av^b-1$ when $N=2^a3^b$.

Let
\begin{equation}\label{eq:Lambda}
 \Lambda=\{N-1:N\in\Ssmooth,\ N\ge2\}.
\end{equation}
This exponent set is too thin for M\"untz completeness.

\begin{lemma}[Convergent reciprocal sum]\label{lem:reciprocal-sum}
The set $\Lambda$ satisfies
\begin{equation}\label{eq:sum-convergent}
 \sum_{m\in\Lambda}\frac1m<\infty.
\end{equation}
More explicitly, the sum is at most $4$.
\end{lemma}

\begin{proof}
For $N\ge2$, $N-1\ge N/2$, so
\[
 \sum_{N\in\Ssmooth,\,N\ge2}\frac1{N-1}
 \le2\sum_{\substack{a,b\ge0\\(a,b)\ne(0,0)}}\frac1{2^a3^b}
 =2\left(\frac1{1-1/2}\frac1{1-1/3}-1\right)=4.
\]
Unique factorization ensures that the parametrization $N=2^a3^b$ has no repetitions.
\end{proof}

\subsection{Why a density theorem is the right audit}

The classical M\"untz--Sz\'asz theorem says, in the present integer-exponent setting, that
\[
 \Span\{1,q^m:m\in\Lambda\}
\]
is dense in $C[0,1]$ precisely when $\sum_{m\in\Lambda}1/m$ diverges.  Full versions give the
same convergent-sum obstruction in $L^p$ and on intervals away from the origin
\cite{Almira2007,BorweinErdelyi1995}.  Lemma~\ref{lem:reciprocal-sum} therefore implies that
the $3$-smooth monomials leave an infinite-dimensional blind subspace.

The kernels $K_m$ differ from $1-q^m$ only by multiplication by $(1-q)^{-1}$.  On any compact
subinterval of $(0,1)$ this multiplier is bounded and invertible.  The same incompleteness
therefore holds for the complete observation family in \eqref{eq:Km-data}.

\section{M\"untz invisibility: an infinite-dimensional no-go theorem}

The next theorem is the strongest conclusion of this report.  It says that even \emph{all}
linear $3$-smooth observations, together with the exact endpoint behavior of the measure, do
not identify $\mu_\alpha$ among positive measures.

\begin{theorem}[Bounded invisible perturbations]\label{thm:muntz-invisible}
Fix $0<\alpha<1$ and a closed interval $I=[a,b]\subset(0,1)$.  There exists a nonzero bounded
real function $g\in L^\infty(\mu_\alpha)$ supported in $I$ such that
\begin{equation}\label{eq:annihilate}
 \int_0^1K_m(q)g(q)\dd\mu_\alpha(q)=0
 \qquad\text{for every }m\in\Lambda.
\end{equation}
Consequently, for every sufficiently small $\varepsilon>0$, the two measures
\begin{equation}\label{eq:mu-pm}
 \dd\mu_\alpha^\pm=(1\pm\varepsilon g)\dd\mu_\alpha
\end{equation}
are distinct positive measures satisfying
\begin{equation}\label{eq:same-data}
 \int_0^1K_m\dd\mu_\alpha^+
 =\int_0^1K_m\dd\mu_\alpha^-
 =\int_0^1K_m\dd\mu_\alpha
 \qquad(m\in\Lambda).
\end{equation}
They agree with $\mu_\alpha$ outside $I$.
\end{theorem}

\begin{proof}
On $I$, multiplication by $1-q$ is a bounded invertible map on $L^1(I)$.  If
$\Span\{K_m:m\in\Lambda\}$ were dense in $L^1(I)$, then
$\Span\{1-q^m:m\in\Lambda\}$ would be dense as well, and adding constants would make
$\Span\{1,q^m:m\in\Lambda\}$ dense.  This contradicts the convergent-sum part of the full
M\"untz theorem and Lemma~\ref{lem:reciprocal-sum}.  Hence the $K_m$ span is a proper
subspace of $L^1(I)$.

By Hahn--Banach and the identification $(L^1(I))^*=L^\infty(I)$, there is a nonzero real
$h\in L^\infty(I)$ such that
\begin{equation}\label{eq:h-ann}
 \int_IK_m(q)h(q)\dd q=0
 \qquad(m\in\Lambda).
\end{equation}
The density $w_\alpha$ in \eqref{eq:mu-density} is continuous and bounded above and below by
positive constants on $I$.  Define
\[
 g(q)=\begin{cases}
 h(q)/w_\alpha(q),&q\in I,\\
 0,&q\notin I.
 \end{cases}
\]
Then $g\in L^\infty(\mu_\alpha)$, is nonzero, and \eqref{eq:annihilate} follows from
\eqref{eq:h-ann}.  Choose $\varepsilon\|g\|_\infty<1$.  Then
$1\pm\varepsilon g>0$ almost everywhere, so \eqref{eq:mu-pm} defines distinct positive
measures.  Equation \eqref{eq:same-data} is immediate from \eqref{eq:annihilate}, and the
support condition gives agreement outside $I$.
\end{proof}

\begin{corollary}[Endpoint asymptotics are also invisible]\label{cor:endpoint-invisible}
The perturbations in Theorem~\ref{thm:muntz-invisible} can be chosen to preserve the density of
$\mu_\alpha$ on neighborhoods of both $0$ and $1$.  Their ordinary moments differ from those
of $\mu_\alpha$ by $O(b^j)$, while
\[
 \int q^j\dd\mu_\alpha(q)\sim\alpha j^{\alpha-1}.
\]
Thus they have the same polynomial moment asymptotics, to an error exponentially smaller than
the main term, and still agree on every $3$-smooth observation \eqref{eq:Km-data}.
\end{corollary}

\begin{proof}
Choose $I=[a,b]$ with $0<a<b<1$.  The measures coincide off $I$.  Moreover
\[
 \left|\int_0^1q^j\dd(\mu_\alpha^+-\mu_\alpha^-)(q)\right|
 \le 2\varepsilon\|g\|_\infty\mu_\alpha(I)b^j.
\]
Combine this with \eqref{eq:tail}.
\end{proof}

\begin{corollary}[Linear measure reconstruction cannot prove \AEname{}]\label{cor:no-linear-proof}
No argument whose only inputs are
\begin{enumerate}[label=\textup{(\alph*)}]
\item positivity of an unknown measure on $[0,1]$;
\item the complete set of scalar observations \eqref{eq:Km-data}; and
\item the local endpoint form or polynomial tail asymptotics of $\mu_\alpha$
\end{enumerate}
can uniquely recover the explicit Hausdorff measure.  In particular, linear functional
analysis, positive kernels, and M\"untz-style measure determination alone cannot settle the
\AEname{} conjecture.
\end{corollary}

\begin{remark}
The theorem does not create another exponent $\beta$ with the same rational powers.  The
perturbed measures need not belong to the one-parameter family $\{\mu_\beta\}$.  Its force is
methodological and exact: the missing proof must exploit the \emph{parametric formula} for
$w_\alpha$ and the arithmetic of its values nonlinearly, because all linear observation data
can be held fixed under positive perturbation.
\end{remark}

\section{An analytic reformulation of the invisible directions}

For a bounded perturbation $g$ supported in $[a,b]\subset(0,1)$, define its Mellin transform
relative to $\mu_\alpha$ by
\begin{equation}\label{eq:G-def}
 G(z)=\int_0^1q^zg(q)\dd\mu_\alpha(q).
\end{equation}
Because the support stays away from $0$, $G$ is an entire function of $z$ and decays
exponentially along the positive real axis.  Condition \eqref{eq:annihilate} becomes
\begin{equation}\label{eq:partial-zero}
 \sum_{j=0}^{m-1}G(j)=0
 \qquad(m\in\Lambda).
\end{equation}
Thus M\"untz incompleteness produces nonzero entire Mellin transforms whose discrete partial
sums vanish at every $3$-smooth-minus-one index.  The zero set is sparse enough to permit this
without forcing $G=0$.

By contrast, the unperturbed transform is explicit:
\[
 M_\alpha(z)=\int_0^1q^z\dd\mu_\alpha(q)
 =(z+2)^\alpha-(z+1)^\alpha.
\]
A future proof would have to combine this special fractional-power form with the arithmetic
statement that
\[
 1+\sum_{j=0}^{N-2}M_\alpha(j)=N^\alpha=u^av^b
 \quad(N=2^a3^b)
\]
has prime-separated rational denominators.  Neither the zero set in
\eqref{eq:partial-zero} nor positivity alone detects that nonlinear coupling.

\section{What would constitute a proof along this route?}

\subsection{The exact nonlinear target}

The moment representation converts the conjecture into the following statement.

\begin{target}[Arithmetic Hausdorff rigidity]\label{target:arithmetic-hausdorff}
Prove that no $\alpha\in(0,1)$ can simultaneously satisfy
\begin{equation}\label{eq:target-coordinates}
 2^\alpha\in\Z[1/2],
 \qquad
 3^\alpha\in\Z[1/3],
\end{equation}
using a property of the explicit density $w_\alpha$ that is nonlinear in the observation
values
\[
 \int K_{2^a3^b-1}\dd\mu_\alpha=2^{a\alpha}3^{b\alpha}-1.
\]
A sufficient theorem must remain sensitive to the separate $2$-primary and $3$-primary
reduced denominators and cannot depend only on linear combinations of the kernels $K_m$.
\end{target}

The target is equivalent to the original problem after the fractional normalization, but the
report's no-go theorems sharply restrict plausible proofs.  A successful invariant could take
one of the following forms:

\begin{enumerate}
\item a nonlinear determinant or resultant built from several \emph{reduced} observations,
      with a lower bound from coprimality that exceeds a sharp analytic upper bound;
\item a rigidity theorem for the explicit Mellin family
      $(z+2)^\alpha-(z+1)^\alpha$ under coherent dyadic/triadic rational partial sums;
\item a mixed real--nonarchimedean identity that transports the exponent $\alpha$ itself,
      rather than merely the rational values, to $2$-adic and $3$-adic analytic data;
\item an unconditional logarithm-product independence theorem strong enough to rule out
      \eqref{eq:log-det}.
\end{enumerate}

The first three are genuinely adapted to the additive moment formulation.  The fourth returns
to the classical four-exponentials frontier.

\subsection{Why the current exact defects do not close}

The new defects have three desirable features: positivity, rationality, and exact
prime-power denominators.  What they lack is an upper bound small enough to force their
positive integer numerators to vanish.  The natural analytic upper bounds deteriorate only
polynomially in the endpoint distance, while the denominator lower bounds are exponentially
small in $r+s$.  Lemma~\ref{lem:trivial-gaps} already supplies comparable or stronger
one-coordinate bounds.  Increasing the finite order cannot change this diagnosis because
Corollary~\ref{cor:finite-simulation} constructs rational atomic models for every finite
collection.

The infinite collection also fails as a linear uniqueness set by
Theorem~\ref{thm:muntz-invisible}.  Therefore the missing gain cannot be ``more positivity''
or ``more moments'' in the same linear sense.  It must be a nonlinear arithmetic identity
that has no analogue for the rational quadrature and invisible perturbation models.

\subsection{A precise conditional resolution}

For completeness, the conjecture does have an immediate conditional proof.

\begin{theorem}[Conditional resolution]\label{thm:conditional-resolution}
The \AEname{} conjecture follows from the Four Exponentials Conjecture, and hence from Schanuel's
conjecture.
\end{theorem}

\begin{proof}
This is Proposition~\ref{prop:minimal-reductions}(iv).  Schanuel's conjecture implies Four
Exponentials by the standard transcendence-degree argument.
\end{proof}

This conditional theorem is not a substitute for an unconditional proof, but it identifies the
known transcendence strength sufficient to close every remaining case.

\section{Comparison with the supplied report}

The contribution of this continuation can be summarized without restating the existing
unified report.

\begin{longtable}{@{}>{\raggedright\arraybackslash}p{0.25\textwidth}>{\raggedright\arraybackslash}p{0.33\textwidth}>{\raggedright\arraybackslash}p{0.34\textwidth}@{}}
\caption{What is inherited, what is new, and what remains.}\label{tab:comparison}\\
\toprule
Topic & Supplied unified report & This continuation\\
\midrule
\endfirsthead
\toprule
Topic & Supplied unified report & This continuation\\
\midrule
\endhead
Four-exponentials barrier
& Exhaustive logarithmic, matrix, and determinant formulations.
& Only the minimal standalone reduction is recalled.\\[0.5em]
Conditional counterexample structure
& Exact free-monoid and valuation classifications.
& Not repeated; fractional normalization is used only to expose $u$ and $v$.\\[0.5em]
Additive relation $3=2+1$
& Appears indirectly in finite differences and Loewner-type methods.
& Converted into an explicit Hausdorff moment measure for all consecutive increments.\\[0.5em]
Adjacent $3$-smooth inputs
& No Hausdorff-moment classification.
& Exact list $(1,2),(2,3),(3,4),(8,9)$ and rational moments $d_0,d_1,d_2,d_7$.\\[0.5em]
Arithmetic positivity
& Many determinant and valuation certificates.
& Exact reduced denominators for one Hankel and three nonconsecutive log-convex defects.\\[0.5em]
Finite positivity barrier
& General determinant-budget obstructions.
& Explicit rational two-atom model and general rational finite quadrature theorem.\\[0.5em]
Infinite linear barrier
& Polynomial approximation and spectral sparsity analyses.
& M\"untz incompleteness gives bounded positive perturbations preserving all $3$-smooth
  observations and endpoint asymptotics.\\[0.5em]
Final status
& Conjecture open; remaining wall isolated.
& Conjecture still open; the moment route is narrowed to nonlinear prime-separated arithmetic.\\
\bottomrule
\end{longtable}

\section{Conclusion}

The requested unconditional implication has not been proved.  The investigation nevertheless
produces a coherent new theory around the additive structure that the standard logarithmic
matrix forgets.

A hypothetical noninteger solution gives a fractional exponent $0<\alpha<1$ for which the
increments of $t^\alpha$ form an explicit Hausdorff moment sequence.  The only consecutive
$3$-smooth pairs isolate four rational moments, and those moments generate positive Hankel and
log-convexity defects whose reduced denominators can be computed exactly; their numerators are positive and coprime to the displayed denominators (and to $6$ when $s>0$).  These are rigorous new arithmetic certificates.

The same analysis also explains why the certificates do not constitute a proof.  The first
three moments possess an explicit rational two-atom model.  Every finite collection of rational
polynomial observations has a rational positive quadrature model.  Most decisively, the
$3$-smooth-minus-one exponents have a convergent reciprocal sum, so the full M\"untz theorem
creates bounded perturbations of the actual measure that preserve all $3$-smooth linear data,
remain positive, and leave the endpoint singularities and polynomial moment tail unchanged.

The remaining problem is therefore not a generic moment problem.  It is a nonlinear arithmetic
rigidity problem for the special one-parameter density $w_\alpha$, with coherent multiplicative
values and mutually separated dyadic and triadic denominators.  Four Exponentials would settle
it immediately.  Without that conjecture, a new theorem of precisely this nonlinear type is
required.

\appendix

\section{Algebraic audit of the exact denominator formulas}

For convenient verification, the common-denominator calculations are collected here.
Starting from \eqref{eq:dA},
\begin{align*}
 d_0d_2-d_1^2
 &=\frac{A_0A_2}{2^{3r}3^s}
   -\frac{A_1^2}{2^{2r}3^{2s}}\\
 &=\frac{3^sA_0A_2-2^rA_1^2}{2^{3r}3^{2s}},\\[0.4em]
 d_0^6d_7-d_1^7
 &=\frac{A_0^6A_7}{2^{9r}3^{2s}}
   -\frac{A_1^7}{2^{7r}3^{7s}}\\
 &=\frac{3^{5s}A_0^6A_7-2^{2r}A_1^7}{2^{9r}3^{7s}},\\[0.4em]
 d_1^5d_7-d_2^6
 &=\frac{A_1^5A_7}{2^{8r}3^{7s}}
   -\frac{A_2^6}{2^{12r}3^{6s}}\\
 &=\frac{2^{4r}A_1^5A_7-3^sA_2^6}{2^{12r}3^{7s}},\\[0.4em]
 d_0^5d_7^2-d_2^7
 &=\frac{A_0^5A_7^2}{2^{11r}3^{4s}}
   -\frac{A_2^7}{2^{14r}3^{7s}}\\
 &=\frac{2^{3r}3^{3s}A_0^5A_7^2-A_2^7}{2^{14r}3^{7s}}.
\end{align*}
The parity and residue arguments in Theorems~\ref{thm:H-denom} and
\ref{thm:higher-defects} show that these common denominators are already reduced.

\section{A moment-cone interpretation}

For a finite set of exponents $E=\{e_1,\ldots,e_m\}$, define the moment curve
\[
 \Phi_E(q)=(q^{e_1},\ldots,q^{e_m})\in\R^m.
\]
The moment vector of $\mu_\alpha$ is
\[
 y_E=\int_0^1\Phi_E(q)\dd\mu_\alpha(q).
\]
Because the density is positive, $y_E$ lies in the interior of the cone generated by
$\Phi_E((0,1))$.  This is the geometric reason rational quadrature works: when $y_E$ happens
to be rational, rational points of the moment curve can be perturbed to surround it, and the
resulting positive barycentric coordinates can be chosen rational.  The transcendental
parameter $\alpha$ is therefore invisible to every finite-dimensional moment cone.

For $E=\{0,1,2\}$, the cone is described by the truncated Hausdorff conditions
\[
 m_0\ge m_1\ge m_2\ge0,
 \qquad m_0m_2\ge m_1^2.
\]
The sequence from Theorem~\ref{thm:hausdorff} lies strictly in the interior.  Formula
\eqref{eq:two-atom-measure} is an explicit rational point of the fiber over the same moment
vector.

\section{Classical inputs and bibliography notes}

The report uses three deep external results.

\begin{enumerate}
\item \textbf{Gelfond--Schneider.}  If $a$ and $b$ are algebraic, $a\ne0,1$, and $b$ is
      algebraic irrational, then every value of $a^b$ is transcendental.  It rules out
      algebraic irrational counterexamples.
\item \textbf{Catalan--Mih\u{a}ilescu.}  The only solution in integers $X,Y,m,n>1$ of
      $X^m-Y^n=1$ is $3^2-2^3=1$.  It classifies the consecutive $3$-smooth pairs.
\item \textbf{Full M\"untz theorem.}  For a sequence of distinct positive exponents bounded
      away from $0$, convergence of the reciprocal sum prevents density of the associated
      monomials in $C$ and $L^p$ spaces, including intervals away from the origin.  It produces
      the invisible directions in Theorem~\ref{thm:muntz-invisible}.
\end{enumerate}

No unproved transcendence conjecture is used in the unconditional theorems.  Four Exponentials
and Schanuel appear only in the explicitly conditional Theorem~\ref{thm:conditional-resolution}.

\begin{thebibliography}{99}

\bibitem{AlaogluErdos1944}
L.~Alaoglu and P.~Erd\H{o}s,
\newblock \emph{On highly composite and similar numbers},
\newblock Transactions of the American Mathematical Society \textbf{56} (1944), 448--469.
\newblock DOI: \href{https://doi.org/10.1090/S0002-9947-1944-0011087-2}
{10.1090/S0002-9947-1944-0011087-2}.

\bibitem{Almira2007}
J.~M.~Almira,
\newblock \emph{M\"untz type theorems I},
\newblock Surveys in Approximation Theory \textbf{3} (2007), 152--194;
\newblock arXiv:0710.3570.

\bibitem{Baker1975}
A.~Baker,
\newblock \emph{Transcendental Number Theory},
\newblock Cambridge University Press, 1975.

\bibitem{BorweinErdelyi1995}
P.~Borwein and T.~Erd\'elyi,
\newblock \emph{Polynomials and Polynomial Inequalities},
\newblock Graduate Texts in Mathematics 161, Springer, 1995.

\bibitem{DasguptaKakde2024}
S.~Dasgupta and M.~Kakde,
\newblock \emph{Ranks of matrices of logarithms of algebraic numbers II: The Matrix
Coefficient Conjecture},
\newblock arXiv:2408.08178, 2024; published in Advances in Mathematics \textbf{487} (2026),
Article 110753.

\bibitem{Hausdorff1921a}
F.~Hausdorff,
\newblock \emph{Summationsmethoden und Momentfolgen. I},
\newblock Mathematische Zeitschrift \textbf{9} (1921), 74--109.

\bibitem{Hausdorff1921b}
F.~Hausdorff,
\newblock \emph{Summationsmethoden und Momentfolgen. II},
\newblock Mathematische Zeitschrift \textbf{9} (1921), 280--299.

\bibitem{Mihailescu2004}
P.~Mih\u{a}ilescu,
\newblock \emph{Primary cyclotomic units and a proof of Catalan's conjecture},
\newblock Journal f\"ur die reine und angewandte Mathematik \textbf{572} (2004), 167--195.
\newblock DOI: \href{https://doi.org/10.1515/crll.2004.048}{10.1515/crll.2004.048}.

\bibitem{ShohatTamarkin1943}
J.~A.~Shohat and J.~D.~Tamarkin,
\newblock \emph{The Problem of Moments},
\newblock Mathematical Surveys 1, American Mathematical Society, 1943.

\bibitem{Waldschmidt2000}
M.~Waldschmidt,
\newblock \emph{Diophantine Approximation on Linear Algebraic Groups},
\newblock Grundlehren der mathematischen Wissenschaften 326, Springer, 2000.

\bibitem{Waldschmidt2023}
M.~Waldschmidt,
\newblock \emph{The Four Exponentials Problem and Schanuel's Conjecture},
\newblock in \emph{Mathematics Going Forward}, Lecture Notes in Mathematics 2313,
Springer, 2023, pp.~579--592.
\newblock DOI: \href{https://doi.org/10.1007/978-3-031-12244-6_39}
{10.1007/978-3-031-12244-6\_39}.

\bibitem{Widder1941}
D.~V.~Widder,
\newblock \emph{The Laplace Transform},
\newblock Princeton University Press, 1941.

\end{thebibliography}

\end{document}
